%%% =======================================================================
%%% ISLS Extension Architecture v1.0.0
%%% Manifest Layer, Registry Infrastructure, Capsule Protocol,
%%% Spiral Scheduler, and Execute Mode
%%% Author: Sebastian Klemm
%%% Date: 16 March 2026
%%% Status: Implementation Specification for Claude Code
%%% =======================================================================
\documentclass[11pt,a4paper,twoside]{article}

\usepackage[utf8]{inputenc}
\usepackage{amsmath,amssymb,amsthm,mathtools}
\usepackage{algorithm,algorithmic}
\usepackage{booktabs,longtable,tabularx}
\usepackage{enumitem}
\usepackage{geometry}
\usepackage{hyperref}
\usepackage{cleveref}
\usepackage{xcolor}
\usepackage{fancyhdr}
\usepackage{titlesec}

\geometry{margin=2.5cm,top=3cm,bottom=3cm,headheight=14pt}

%%% Theorem environments
\theoremstyle{definition}
\newtheorem{definition}{Definition}[section]
\newtheorem{axiom}{Axiom}[section]
\newtheorem{requirement}{Requirement}[section]
\newtheorem{invariant}{Invariant}[section]

\theoremstyle{plain}
\newtheorem{theorem}{Theorem}[section]
\newtheorem{proposition}{Proposition}[section]
\newtheorem{lemma}{Lemma}[section]

\theoremstyle{remark}
\newtheorem{remark}{Remark}[section]

%%% Commands
\newcommand{\ISLS}{\textsf{ISLS}}
\newcommand{\RR}{\mathbb{R}}
\newcommand{\NN}{\mathbb{N}}
\newcommand{\ZZ}{\mathbb{Z}}
\newcommand{\Hfive}{H^{5}}
\newcommand{\calO}{\mathcal{O}}
\newcommand{\calR}{\mathcal{R}}
\newcommand{\calM}{\mathcal{M}}
\newcommand{\calC}{\mathcal{C}}
\newcommand{\calK}{\mathcal{K}}
\newcommand{\calP}{\mathcal{P}}
\newcommand{\calS}{\mathcal{S}}
\newcommand{\calE}{\mathcal{E}}
\newcommand{\calL}{\mathcal{L}}
\newcommand{\calG}{\mathcal{G}}
\newcommand{\calA}{\mathcal{A}}
\newcommand{\OmegaISLS}{\Omega_{\ISLS}}
\DeclareMathOperator{\Canon}{Canon}
\DeclareMathOperator{\Dig}{Dig}
\DeclareMathOperator{\HKDF}{HKDF}
\DeclareMathOperator{\AEAD}{AEAD}
\DeclareMathOperator{\Replay}{Replay}
\DeclareMathOperator{\Commit}{Commit}
\newcommand{\norm}[1]{\left\lVert #1 \right\rVert}

%%% Header
\pagestyle{fancy}
\fancyhf{}
\fancyhead[LE,RO]{\ISLS{} Extension Architecture v1.0.0}
\fancyhead[RE,LO]{Sebastian Klemm}
\fancyfoot[C]{\thepage}

\hypersetup{
  colorlinks=true,
  linkcolor=blue!60!black,
  pdftitle={ISLS Extension Architecture v1.0.0},
  pdfauthor={Sebastian Klemm},
}

\begin{document}

%%% Title
\begin{titlepage}
\centering
\vspace*{2.5cm}
{\Huge\bfseries \ISLS{} Extension Architecture\\[0.3cm]
v1.0.0}\\[1.5cm]
{\Large Manifest Layer, Registry Infrastructure,\\
Capsule Protocol, Spiral Scheduler, and Execute Mode}\\[1cm]
{\large A Complementary Architecture for the\\
Intelligent Semantic Ledger Substrate}\\[2cm]
{\large Sebastian Klemm}\\[0.5cm]
{\large 16 March 2026}\\[2cm]

\begin{abstract}
\noindent
This document specifies four extension crates and one engine modification
for the Intelligent Semantic Ledger Substrate (\ISLS{}). Together they
complete the PSP (Post-Symbolic Programming) substrate by adding:
(1)~a run-level Execution Manifest binding all artifacts of a complete run
into a single content-addressed meta-artifact;
(2)~a formal digest-bound Registry infrastructure for operators, profiles,
and obligations;
(3)~a generic evidence-bound Capsule protocol (OLP) enabling
policy-controlled secret release gated by deterministic run evidence;
(4)~a Spiral Scheduler providing adaptive tick granularity;
and (5)~an Execute mode enabling generative program execution alongside
the existing analytical discovery mode.

All extensions preserve the 20 structural invariants (I1--I20), all 130+
acceptance tests, deterministic replay, and append-only semantics of the
\ISLS{} core (C1--C11). The specification is implementation-ready for
autonomous coding agents.

\medskip\noindent\textbf{Status:} Implementation Specification.\\
\textbf{Parent:} ISLS v1.0.0 Canonical Master Specification.\\
\textbf{Target:} Rust workspace under \texttt{isls/crates/}.
\end{abstract}
\end{titlepage}

\tableofcontents
\newpage

%%% =====================================================================
\part{Foundations}\label{part:foundations}

\section{Purpose and Scope}\label{sec:purpose}

\subsection{Purpose}

This specification defines the minimal set of extensions required to
evolve \ISLS{} from a domain-specific analytical engine into a
\emph{universal PSP substrate} capable of three operational
manifestations:

\begin{enumerate}[label=(\alph*)]
  \item \textbf{Analytical} (existing): discover programs latent in
        observed data and crystallize them.
  \item \textbf{Generative} (new): accept an operator graph or crystal
        as input and execute it through the staged closure, producing
        new artifacts.
  \item \textbf{Cryptographic} (new): bind secrets to evidence
        certificates produced by analytical or generative runs, releasing
        them only upon verified policy satisfaction.
\end{enumerate}

All three manifestations share: the same state space $\Sigma_0 = \RR^5$,
the same canonicalization (JCS + SHA-256), the same Run Descriptor, the
same Evidence Chain, the same Crystal structure, and the same gate
cascade.

\subsection{Scope}

\begin{itemize}[nosep]
  \item \textbf{In scope:} Four new crates (\texttt{isls-manifest},
        \texttt{isls-registry}, \texttt{isls-capsule},
        \texttt{isls-scheduler}) and two engine extensions (Spiral
        Scheduler integration, Execute mode).
  \item \textbf{Out of scope:} Language frontends (Sanskroot or other
        surface syntaxes), domain-specific operator implementations,
        cryptographic primitive design, external KMS/HSM integration
        details.
\end{itemize}

\subsection{Relationship to Existing Crates}

\begin{tabularx}{\textwidth}{l l X}
\toprule
\textbf{Existing} & \textbf{New} & \textbf{Relationship} \\
\midrule
\texttt{isls-types} & \texttt{isls-registry} & Registry types extend
  \texttt{isls-types}; all registry entries use existing
  \texttt{content\_address}. \\
\texttt{isls-archive} & \texttt{isls-manifest} & Manifest reads the
  crystal archive and produces a run-level meta-artifact. \\
\texttt{isls-consensus} & \texttt{isls-capsule} & Capsule policies
  reference gate evidence produced by \texttt{isls-consensus}. \\
\texttt{isls-engine} & \texttt{isls-scheduler} & Scheduler wraps the
  macro-step loop; engine gains \texttt{execute} mode. \\
\bottomrule
\end{tabularx}

\subsection{Invariant Preservation}

\begin{requirement}[Extension Non-Regression]\label{req:nonregression}
All extensions \textbf{MUST} preserve invariants I1--I20 of the
\ISLS{} v1.0.0 specification. All existing acceptance tests (AT-01
through AT-20) and all existing unit/integration tests \textbf{MUST}
continue to pass without modification.
\end{requirement}

\subsection{Technical Constraints}

\begin{requirement}[Workspace Discipline]\label{req:workspace}
All new crates \textbf{MUST} adhere to the established workspace
discipline:
\begin{itemize}[nosep]
  \item No \texttt{unsafe} Rust code.
  \item No \texttt{rand} dependency; determinism via \texttt{RunDescriptor}.
  \item \texttt{BTreeMap} for all deterministic-order collections.
  \item JCS (RFC~8785) canonical serialization via \texttt{serde\_jcs}.
  \item SHA-256 content addressing via \texttt{sha2}.
  \item Append-only semantics; no \texttt{delete} methods.
  \item All operators version-pinned.
\end{itemize}
\end{requirement}


%%% =====================================================================
\part{Crate C12: \texttt{isls-registry}}\label{part:registry}

\section{Purpose}\label{sec:reg:purpose}

The registry crate provides a formal, digest-bound catalog system for
operators, profiles, and obligations. It replaces the flat
\texttt{operator\_versions: BTreeMap<String, String>} in existing
crystals with a structured, content-addressed registry infrastructure
that satisfies the PCR registry requirements.

\section{Data Model}\label{sec:reg:model}

\begin{definition}[Registry Entry]\label{def:regentry}
A registry entry is a tuple
\[
  e = \langle \mathtt{id}, \mathtt{name}, \mathtt{version}, \mathtt{digest},
      \mathtt{kind}, \mathtt{metadata} \rangle
\]
where:
\begin{itemize}[nosep]
  \item \texttt{id}: content-address $\Dig(\Canon(e_{\mathrm{core}}))$,
  \item \texttt{name}: human-readable identifier (e.g., \texttt{"BandOperator"}),
  \item \texttt{version}: semantic version string (e.g., \texttt{"1.0.0"}),
  \item \texttt{digest}: SHA-256 of the canonical serialization of the
        entry's operational content,
  \item \texttt{kind} $\in \{\texttt{Operator}, \texttt{Profile},
        \texttt{Obligation}, \texttt{Macro}\}$,
  \item \texttt{metadata}: \texttt{BTreeMap<String, String>} for
        extensible key-value annotations.
\end{itemize}
\end{definition}

\begin{definition}[Registry]\label{def:registry}
A registry $\calR$ is an append-only, content-addressed collection of
registry entries:
\[
  \calR = \langle \mathtt{registry\_id}, \mathtt{kind}, \mathtt{entries},
           \mathtt{digest} \rangle
\]
where $\mathtt{registry\_id} = \Dig(\Canon(\mathtt{entries}))$ and
$\mathtt{digest}$ is the SHA-256 of the entire canonical registry
serialization.
\end{definition}

\begin{definition}[Registry Set]\label{def:regset}
A registry set is a tuple of four registries:
\[
  \calR_{\mathrm{all}} = (\calR_{\mathrm{op}}, \calR_{\mathrm{prof}},
  \calR_{\mathrm{obl}}, \calR_{\mathrm{macro}})
\]
bound in the Run Descriptor by their respective digests.
\end{definition}

\section{Operator Registry}\label{sec:reg:operator}

\begin{definition}[Operator Registration]\label{def:opreg}
Each operator in $\OmegaISLS$ \textbf{MUST} be registered with:
\begin{itemize}[nosep]
  \item \texttt{name}: canonical operator name,
  \item \texttt{version}: semantic version,
  \item \texttt{type\_signature}: input type $\to$ output type,
  \item \texttt{role}: one of $\{\texttt{ingest}, \texttt{canon},
        \texttt{measure}, \texttt{solve}, \texttt{morph}, \texttt{gate},
        \texttt{select}, \texttt{merge}, \texttt{emit}\}$ (PSP kernel
        roles),
  \item \texttt{determinism\_certificate}: boolean assertion that the
        operator is deterministic under fixed $(\Sigma, \mathrm{RD})$,
  \item \texttt{digest}: SHA-256 of the canonical serialization of the
        above fields.
\end{itemize}
\end{definition}

\begin{requirement}[Registry Binding in RD]\label{req:rdbinding}
The Run Descriptor \textbf{MUST} contain the digests of all four
registries. Any operator, profile, obligation, or macro used during a
run \textbf{MUST} be resolvable from the registry referenced in the RD.
An operator not present in the registry \textbf{MUST NOT} be executed.
\end{requirement}

\begin{requirement}[Drift Detection]\label{req:drift}
If an operator's implementation digest does not match its registry entry
digest, the system \textbf{MUST} raise an operator-drift fault
(Invariant~I20). This check \textbf{MUST} occur before operator
execution, not after.
\end{requirement}

\section{Profile Registry}\label{sec:reg:profile}

\begin{definition}[Profile Entry]\label{def:profentry}
A profile entry binds a named configuration to a specific parameter set:
\begin{itemize}[nosep]
  \item \texttt{name}: one of $\{\texttt{Minimal}, \texttt{Standard},
        \texttt{HighRigor}\}$ or custom,
  \item \texttt{thresholds}: all gate threshold values,
  \item \texttt{normalization}: all normalization parameters
        ($\mu_D, \mu_Q, \lambda, \mu_r, \ldots$),
  \item \texttt{carrier}: carrier geometry configuration,
  \item \texttt{consensus}: dual consensus parameters,
  \item \texttt{digest}: SHA-256 of canonical serialization.
\end{itemize}
\end{definition}

\section{Obligation Registry}\label{sec:reg:obligation}

\begin{definition}[Obligation Entry]\label{def:oblentry}
An obligation entry defines a proof requirement:
\begin{itemize}[nosep]
  \item \texttt{name}: obligation identifier (e.g.,
        \texttt{"evidence\_chain\_complete"}),
  \item \texttt{gate\_ids}: which gates must produce evidence,
  \item \texttt{evidence\_schema}: canonical JSON schema of required
        evidence fields,
  \item \texttt{verification\_procedure}: reference to verifier function
        name and version,
  \item \texttt{digest}: SHA-256 of canonical serialization.
\end{itemize}
\end{definition}

\section{Rust API}\label{sec:reg:api}

\begin{verbatim}
// isls-registry/src/lib.rs

pub struct RegistryEntry {
    pub id: Hash256,
    pub name: String,
    pub version: String,
    pub digest: Hash256,
    pub kind: RegistryKind,
    pub metadata: BTreeMap<String, String>,
}

pub enum RegistryKind { Operator, Profile, Obligation, Macro }

pub struct Registry {
    pub registry_id: Hash256,
    pub kind: RegistryKind,
    pub entries: BTreeMap<String, RegistryEntry>,  // name -> entry
    pub digest: Hash256,
}

pub struct RegistrySet {
    pub operators: Registry,
    pub profiles: Registry,
    pub obligations: Registry,
    pub macros: Registry,
}

impl Registry {
    pub fn new(kind: RegistryKind) -> Self;
    pub fn register(&mut self, entry: RegistryEntry) -> Result<()>;
    pub fn resolve(&self, name: &str) -> Option<&RegistryEntry>;
    pub fn verify_digest(&self, name: &str, expected: &Hash256) -> bool;
    pub fn compute_digest(&self) -> Hash256;
}

impl RegistrySet {
    pub fn from_config(config: &Config) -> Self;
    pub fn digests(&self) -> BTreeMap<String, Hash256>;
    pub fn verify_operator(&self, name: &str, digest: &Hash256) -> bool;
}
\end{verbatim}

\section{Acceptance Tests}\label{sec:reg:tests}

\begin{tabularx}{\textwidth}{l l X}
\toprule
\textbf{ID} & \textbf{Name} & \textbf{Procedure} \\
\midrule
AT-R1 & Registry content address & Register an operator; verify
  \texttt{id} = SHA-256(JCS(core fields)). \\
AT-R2 & Drift detection & Modify operator implementation; verify drift
  fault raised before execution. \\
AT-R3 & RD binding & Create RD with registry digests; verify run rejects
  unregistered operator. \\
AT-R4 & Append-only & Attempt to remove registry entry; verify
  rejection. \\
AT-R5 & Deterministic digest & Register same entries in different order;
  verify identical registry digest (BTreeMap ordering). \\
\bottomrule
\end{tabularx}


%%% =====================================================================
\part{Crate C13: \texttt{isls-manifest}}\label{part:manifest}

\section{Purpose}\label{sec:man:purpose}

The manifest crate produces a single content-addressed meta-artifact that
binds an entire run---all crystals, all traces, all evidence, all
registry references---into one verifiable object. This satisfies the PCR
Execution Manifest requirement and enables offline replay without
external queries.

\section{Data Model}\label{sec:man:model}

\begin{definition}[Execution Manifest]\label{def:manifest}
An execution manifest is a tuple
\[
  M = \langle \mathtt{run\_id}, \mathtt{rd\_digest}, \mathtt{program\_id},
      \mathtt{input\_digest}, \mathtt{crystal\_digests},
      \mathtt{trace\_digests}, \mathtt{registry\_digests},
      \mathtt{mef\_head}, \mathtt{timestamp} \rangle
\]
where:
\begin{itemize}[nosep]
  \item $\mathtt{run\_id} = \Dig(\Canon(M_{\mathrm{core}}))$: content
        address of the manifest,
  \item $\mathtt{rd\_digest}$: SHA-256 of the canonical Run Descriptor,
  \item $\mathtt{program\_id}$: identifier of the program executed
        (crystal ID for execute mode, \texttt{"discovery"} for
        analytical mode),
  \item $\mathtt{input\_digest}$: SHA-256 of the canonical observation
        log or input program,
  \item $\mathtt{crystal\_digests}$: ordered list of all crystal IDs
        produced during the run,
  \item $\mathtt{trace\_digests}$: ordered list of per-tick trace
        digests,
  \item $\mathtt{registry\_digests}$: map of registry kind $\to$ digest,
  \item $\mathtt{mef\_head}$: head digest of the MEF evidence chain
        (optional),
  \item $\mathtt{timestamp}$: extrinsic time of manifest creation.
\end{itemize}
\end{definition}

\begin{definition}[Trace Entry]\label{def:trace}
A trace entry records one macro-step:
\[
  T_k = \langle k, \mathtt{input\_digest}_k,
        \mathtt{state\_digest}_k, \mathtt{crystal\_id}_k,
        \mathtt{gate\_snapshot}_k, \mathtt{metrics\_digest}_k \rangle
\]
where $\mathtt{crystal\_id}_k$ is \texttt{None} if no crystal was
committed at step $k$.
\end{definition}

\begin{definition}[Replay Pack]\label{def:replaypack}
A replay pack is a self-contained bundle sufficient for offline replay:
\[
  \mathrm{RP} = \langle M, \mathrm{RD}, \mathtt{observation\_log},
  \calR_{\mathrm{all}}, \mathtt{boundary\_evidence} \rangle
\]
A replay pack \textbf{MUST} enable deterministic reproduction of the run
without any external queries.
\end{definition}

\section{Construction Algorithm}\label{sec:man:algo}

\begin{algorithm}[H]
\caption{\textsc{BuildManifest}: Construct execution manifest post-run}
\label{alg:buildmanifest}
\begin{algorithmic}[1]
\REQUIRE Completed run state $S_K$, Run Descriptor $\mathrm{RD}$,
  trace entries $\{T_0, \ldots, T_{K-1}\}$, archive $A$,
  registry set $\calR_{\mathrm{all}}$
\ENSURE Execution manifest $M$
\STATE $\mathtt{rd\_digest} \gets \Dig(\Canon(\mathrm{RD}))$
\STATE $\mathtt{input\_digest} \gets \Dig(\Canon(\mathtt{observation\_log}))$
\STATE $\mathtt{crystal\_digests} \gets [c.\mathtt{crystal\_id} \mid c \in A.\mathtt{crystals}]$
\STATE $\mathtt{trace\_digests} \gets [\Dig(\Canon(T_k)) \mid k = 0, \ldots, K-1]$
\STATE $\mathtt{registry\_digests} \gets \calR_{\mathrm{all}}.\mathtt{digests}()$
\STATE $\mathtt{mef\_head} \gets A.\mathtt{evidence\_chain\_head}()$
\STATE $M_{\mathrm{core}} \gets (\mathtt{rd\_digest}, \mathtt{input\_digest},
  \mathtt{crystal\_digests}, \mathtt{trace\_digests},
  \mathtt{registry\_digests}, \mathtt{mef\_head})$
\STATE $\mathtt{run\_id} \gets \Dig(\Canon(M_{\mathrm{core}}))$
\STATE $M \gets (\mathtt{run\_id}, M_{\mathrm{core}}, \mathtt{timestamp})$
\RETURN $M$
\end{algorithmic}
\end{algorithm}

\section{Verification}\label{sec:man:verify}

\begin{definition}[Manifest Verification]\label{def:manverify}
A manifest verifier \textbf{MUST} check, in order:
\begin{enumerate}[nosep,label=(MV\arabic*)]
  \item \texttt{run\_id} equals $\Dig(\Canon(M_{\mathrm{core}}))$.
  \item \texttt{rd\_digest} matches the provided RD.
  \item Every entry in \texttt{crystal\_digests} resolves in the archive.
  \item Every trace digest is recomputable from the corresponding trace
        entry.
  \item Registry digests match the provided registries.
  \item If \texttt{mef\_head} is present, the evidence chain is
        verifiable from head to genesis.
\end{enumerate}
On first divergence, the verifier \textbf{MUST} report which check
failed and at which index.
\end{definition}

\section{Rust API}\label{sec:man:api}

\begin{verbatim}
// isls-manifest/src/lib.rs

pub struct ExecutionManifest {
    pub run_id: Hash256,
    pub rd_digest: Hash256,
    pub program_id: String,
    pub input_digest: Hash256,
    pub crystal_digests: Vec<Hash256>,
    pub trace_digests: Vec<Hash256>,
    pub registry_digests: BTreeMap<String, Hash256>,
    pub mef_head: Option<Hash256>,
    pub timestamp: f64,
}

pub struct TraceEntry {
    pub tick: u64,
    pub input_digest: Hash256,
    pub state_digest: Hash256,
    pub crystal_id: Option<Hash256>,
    pub gate_snapshot: GateSnapshot,
    pub metrics_digest: Hash256,
}

pub struct ReplayPack {
    pub manifest: ExecutionManifest,
    pub rd: RunDescriptor,
    pub observation_log: Vec<Vec<Vec<u8>>>,
    pub registries: RegistrySet,
    pub boundary_evidence: Vec<EvidenceEntry>,
}

pub fn build_manifest(
    rd: &RunDescriptor,
    traces: &[TraceEntry],
    archive: &Archive,
    registries: &RegistrySet,
    program_id: &str,
    observation_log: &[Vec<Vec<u8>>],
) -> ExecutionManifest;

pub fn verify_manifest(
    manifest: &ExecutionManifest,
    rd: &RunDescriptor,
    archive: &Archive,
    traces: &[TraceEntry],
    registries: &RegistrySet,
) -> Result<(), ManifestError>;

pub fn build_replay_pack(
    manifest: ExecutionManifest,
    rd: RunDescriptor,
    observation_log: Vec<Vec<Vec<u8>>>,
    registries: RegistrySet,
    boundary_evidence: Vec<EvidenceEntry>,
) -> ReplayPack;
\end{verbatim}

\section{Acceptance Tests}\label{sec:man:tests}

\begin{tabularx}{\textwidth}{l l X}
\toprule
\textbf{ID} & \textbf{Name} & \textbf{Procedure} \\
\midrule
AT-M1 & Manifest content address & Build manifest; verify
  \texttt{run\_id} = SHA-256(JCS(core)). \\
AT-M2 & Manifest verification & Build and verify manifest; all MV1--MV6
  pass. \\
AT-M3 & Tamper detection & Alter one crystal digest; verify
  verification fails at MV3. \\
AT-M4 & Replay pack sufficiency & Build replay pack; replay from pack
  alone; verify identical crystals. \\
AT-M5 & Trace determinism & Run twice with same RD; verify identical
  trace digests. \\
\bottomrule
\end{tabularx}


%%% =====================================================================
\part{Crate C14: \texttt{isls-capsule}}\label{part:capsule}

\section{Purpose}\label{sec:cap:purpose}

The capsule crate implements the Operator Lock Protocol (OLP): a generic
evidence-bound access control mechanism in which secrets are released
only upon verified satisfaction of a deterministic program's gate
obligations. The capsule is domain-agnostic; any \ISLS{} program
(analytical or generative) can serve as a lock.

\section{Threat Model}\label{sec:cap:threat}

\begin{definition}[Adversary Capabilities]\label{def:adversary}
An adversary may:
\begin{enumerate}[nosep,label=(\roman*)]
  \item read the lock program $P_L$, capsule headers, policies, and
        manifests,
  \item copy, replay, and attempt to tamper with bundles and capsules,
  \item run the verifier and the lock program arbitrarily often,
  \item attempt offline attacks on ciphertexts.
\end{enumerate}
\end{definition}

\begin{definition}[Security Objectives]\label{def:secobjectives}
OLP provides:
\begin{enumerate}[nosep,label=(\roman*)]
  \item \textbf{Confidentiality}: secrets remain confidential without
        access to cryptographic key material.
  \item \textbf{Policy-bound release}: \texttt{open} succeeds only if
        lock obligations pass and evidence verifies.
  \item \textbf{Tamper evidence}: any modification to capsule, bundle,
        or manifest is detected.
  \item \textbf{Replay stability}: verification results are invariant
        under replay (deterministic).
\end{enumerate}
\end{definition}

\section{Core Objects}\label{sec:cap:objects}

\begin{definition}[Lock]\label{def:lock}
A lock is a tuple
\[
  L = (P_L, \Gamma, \mathrm{RD}_L)
\]
where $P_L$ is a public operator program (constraint program or HDAG),
$\Gamma$ is the set of gate obligations and thresholds, and
$\mathrm{RD}_L$ is the RD profile binding canonicalization, hash
algorithms, tie-break rules, and registry digests.
\end{definition}

\begin{definition}[Capsule]\label{def:capsule}
A capsule is an encrypted artifact:
\[
  C = \langle \mathtt{schema\_version}, \mathtt{alg}, \mathtt{nonce},
      \mathtt{ciphertext}, \mathtt{aad}, \mathtt{policy},
      \mathtt{bind} \rangle
\]
where:
\begin{itemize}[nosep]
  \item \texttt{alg}: AEAD algorithm identifier (default:
        \texttt{AES-256-GCM}),
  \item \texttt{nonce}: unique per capsule,
  \item \texttt{ciphertext}: $\AEAD.\mathrm{Enc}(K_S, \mathtt{nonce},
        \mathtt{plaintext}, \mathtt{aad})$,
  \item \texttt{aad}: canonical JSON of policy and bind fields
        (authenticated but not encrypted),
  \item \texttt{policy}: capsule policy $\Pi_C$,
  \item \texttt{bind}: digests binding the capsule to PSP objects.
\end{itemize}
\end{definition}

\begin{definition}[Capsule Policy]\label{def:policy}
A capsule policy $\Pi_C$ constrains admissible opening contexts:
\begin{itemize}[nosep]
  \item \texttt{require\_lock\_program\_id}: digest of $P_L$,
  \item \texttt{require\_rd\_digest}: digest of $\mathrm{RD}_L$,
  \item \texttt{require\_gate\_proofs}: list of required gate evidence
        item digests,
  \item \texttt{require\_manifest\_id}: optional; digest of required
        execution manifest,
  \item \texttt{expires\_at}: optional; UTC timestamp after which opening
        is forbidden,
  \item \texttt{max\_uses}: optional; maximum number of successful opens.
\end{itemize}
\end{definition}

\section{Key Schedule}\label{sec:cap:keys}

\begin{definition}[Session Key Derivation]\label{def:keysched}
Given a master key $K_M$ (external, e.g., KMS/HSM/local keychain), define
the binding string:
\[
  B = \Dig(\mathrm{RD}_L) \;\|\; \mathtt{run\_id} \;\|\;
      P_L^{\mathrm{id}} \;\|\; T^{\mathrm{dig}} \;\|\;
      \mathrm{RP}^{\mathrm{dig}}
\]
where values are taken from the execution manifest. Then derive:
\[
  K_S = \HKDF(K_M, \; \mathrm{salt} = \Dig(\Canon(\Sigma_0)), \;
        \mathrm{info} = B)
\]
\end{definition}

\begin{requirement}[Binding Minimality]\label{req:bindmin}
$B$ \textbf{MUST} include at least $\Dig(\mathrm{RD}_L)$ and
$\mathtt{run\_id}$. If boundary operators are used, the replay pack
digest \textbf{MUST} be included to prevent opening under missing
boundary evidence.
\end{requirement}

\section{Protocol Procedures}\label{sec:cap:procedures}

\subsection{Seal}

\begin{algorithm}[H]
\caption{\textsc{Seal}: Encrypt secret under evidence-bound policy}
\label{alg:seal}
\begin{algorithmic}[1]
\REQUIRE Secret $S$, lock $L = (P_L, \Gamma, \mathrm{RD}_L)$, master key $K_M$
\ENSURE Capsule $C$
\STATE Execute $P_L$ on $\Sigma_0$ under $\mathrm{RD}_L$, producing
  manifest $M$, trace $T$, evidence, and replay pack $\mathrm{RP}$
\STATE Verify all gate obligations $\Gamma$ pass with evidence
\STATE Derive $K_S$ from $K_M$ and bindings $B$ (\cref{def:keysched})
\STATE $\mathtt{nonce} \gets$ deterministic nonce from
  $\Dig(\mathtt{run\_id} \| \mathtt{seal\_counter})$
\STATE $\mathtt{aad} \gets \Canon(\Pi_C \| \mathtt{bind})$
\STATE $\mathtt{ciphertext} \gets \AEAD.\mathrm{Enc}(K_S, \mathtt{nonce}, S, \mathtt{aad})$
\STATE $C \gets (\mathtt{schema\_version}, \mathtt{alg}, \mathtt{nonce},
  \mathtt{ciphertext}, \mathtt{aad}, \Pi_C, \mathtt{bind})$
\RETURN $C$
\end{algorithmic}
\end{algorithm}

\subsection{Open}

\begin{algorithm}[H]
\caption{\textsc{Open}: Decrypt secret under verified evidence}
\label{alg:open}
\begin{algorithmic}[1]
\REQUIRE Capsule $C$, lock $L$, manifest $M$, master key $K_M$
\ENSURE Secret $S$ or $\bot$
\STATE Verify manifest $M$ via MV1--MV6 (\cref{def:manverify})
\STATE Check policy $\Pi_C$: program ID, RD digest, gate proofs
\STATE If \texttt{expires\_at} set and current time exceeds it:
  \RETURN $\bot$
\STATE If \texttt{max\_uses} set and exceeded: \RETURN $\bot$
\STATE Verify boundary replay sufficiency if boundary IO was used
\STATE Derive $K_S$ using same derivation rule (\cref{def:keysched})
\STATE $S \gets \AEAD.\mathrm{Dec}(K_S, \mathtt{nonce}, \mathtt{ciphertext}, \mathtt{aad})$
\STATE If decryption fails: \RETURN $\bot$
\RETURN $S$
\end{algorithmic}
\end{algorithm}

\begin{invariant}[Capsule Integrity]\label{inv:capsule}
Any modification to the capsule's \texttt{aad}, \texttt{policy}, or
\texttt{bind} fields \textbf{MUST} cause $\AEAD.\mathrm{Dec}$ to fail,
because these fields are authenticated as associated data.
\end{invariant}

\section{Rust API}\label{sec:cap:api}

\begin{verbatim}
// isls-capsule/src/lib.rs

pub struct Lock {
    pub program: ConstraintProgram,  // or HDAG
    pub obligations: Vec<String>,    // obligation registry names
    pub rd: RunDescriptor,
}

pub struct CapsulePolicy {
    pub require_lock_program_id: Hash256,
    pub require_rd_digest: Hash256,
    pub require_gate_proofs: Vec<Hash256>,
    pub require_manifest_id: Option<Hash256>,
    pub expires_at: Option<f64>,
    pub max_uses: Option<u64>,
}

pub struct Capsule {
    pub schema_version: String,
    pub alg: String,
    pub nonce: Vec<u8>,
    pub ciphertext: Vec<u8>,
    pub aad: Vec<u8>,
    pub policy: CapsulePolicy,
    pub bind: BTreeMap<String, Hash256>,
}

pub fn seal(
    secret: &[u8],
    lock: &Lock,
    master_key: &[u8; 32],
    manifest: &ExecutionManifest,
) -> Result<Capsule, CapsuleError>;

pub fn open(
    capsule: &Capsule,
    lock: &Lock,
    master_key: &[u8; 32],
    manifest: &ExecutionManifest,
) -> Result<Vec<u8>, CapsuleError>;

pub fn verify_policy(
    capsule: &Capsule,
    manifest: &ExecutionManifest,
) -> Result<(), PolicyError>;
\end{verbatim}

\section{Acceptance Tests}\label{sec:cap:tests}

\begin{tabularx}{\textwidth}{l l X}
\toprule
\textbf{ID} & \textbf{Name} & \textbf{Procedure} \\
\midrule
AT-C1 & Seal-open round-trip & Seal a secret; open with valid manifest;
  verify identical plaintext. \\
AT-C2 & Policy rejection & Open with manifest missing a required gate
  proof; verify rejection. \\
AT-C3 & Tamper detection & Alter capsule AAD; verify AEAD decryption
  fails. \\
AT-C4 & Expiry enforcement & Set \texttt{expires\_at} in the past; verify
  open returns $\bot$. \\
AT-C5 & Replay stability & Seal and open twice with same inputs; verify
  identical results. \\
AT-C6 & Wrong manifest & Open with manifest from a different run; verify
  rejection. \\
\bottomrule
\end{tabularx}


%%% =====================================================================
\part{Crate C15: \texttt{isls-scheduler}}\label{part:scheduler}

\section{Purpose}\label{sec:sched:purpose}

The scheduler crate provides adaptive tick granularity based on system
dynamics. Instead of flat ticks of uniform duration, the Spiral Scheduler
observes metric behavior and adjusts tick frequency: finer granularity
during high deformation/friction (turbulence), coarser during stability.

\section{Formal Model}\label{sec:sched:model}

\begin{definition}[Spiral Segmentation]\label{def:spiralseg}
A spiral segmentation maps macro-step index $k$ to an effective tick
count $n_k \in \NN_{>0}$ through a scheduling function:
\[
  n_k = \sigma(d_k, F^*_k, S^*_k, \theta_{\mathrm{sched}})
\]
where $d_k$ is the deformation metric, $F^*_k$ is friction, $S^*_k$ is
shock, and $\theta_{\mathrm{sched}}$ collects scheduling parameters.
\end{definition}

\begin{definition}[Scheduling Function]\label{def:schedfunc}
The canonical scheduling function is:
\[
  \sigma(d, F, S, \theta) = \mathrm{clamp}\!\left(
    \left\lfloor n_{\min} + (n_{\max} - n_{\min}) \cdot
    \max(d, F, S) \right\rfloor,\;
    n_{\min},\; n_{\max}
  \right)
\]
where $n_{\min} \ge 1$ is the minimum sub-steps per macro-step and
$n_{\max}$ is the maximum.
\end{definition}

\begin{definition}[Sub-Step]\label{def:substep}
A sub-step executes the intrinsic dynamics
$\frac{dH}{dt_2} = \Phi(H)$
with effective $\Delta t_2' = \Delta t_2 / n_k$, updating carrier phase,
embeddings, and resonance state $n_k$ times before the extrinsic
commit decision is evaluated.
\end{definition}

\begin{invariant}[Extrinsic Discreteness]\label{inv:extrinsic}
The extrinsic commit index $k$ increments exactly once per macro-step
regardless of $n_k$. Sub-steps are intrinsic-time refinements only.
They \textbf{MUST NOT} produce commits, crystals, or append to the
archive. Only the final composite state after all $n_k$ sub-steps enters
the extrinsic decision pipeline (extraction $\to$ consensus $\to$
commit).
\end{invariant}

\begin{invariant}[Deterministic Scheduling]\label{inv:detsched}
The scheduling function $\sigma$ \textbf{MUST} be deterministic under
fixed $(\Sigma, \mathrm{RD})$. The scheduling parameters
$\theta_{\mathrm{sched}}$ \textbf{MUST} be part of the Run Descriptor.
\end{invariant}

\section{Configuration}\label{sec:sched:config}

\begin{verbatim}
scheduler:
  enabled: true
  n_min: 1
  n_max: 10
  strategy: "max_pressure"  # or "weighted", "fixed"
\end{verbatim}

When \texttt{enabled: false}, the scheduler degenerates to
$n_k = 1$ for all $k$ (flat ticks, backward-compatible).

\section{Rust API}\label{sec:sched:api}

\begin{verbatim}
// isls-scheduler/src/lib.rs

pub struct SchedulerConfig {
    pub enabled: bool,
    pub n_min: u32,
    pub n_max: u32,
    pub strategy: ScheduleStrategy,
}

pub enum ScheduleStrategy {
    MaxPressure,   // n_k = clamp(n_min + (n_max-n_min)*max(d,F,S))
    Weighted { w_d: f64, w_f: f64, w_s: f64 },
    Fixed(u32),    // constant sub-steps
}

pub fn compute_substeps(
    d: f64, f_friction: f64, s_shock: f64,
    config: &SchedulerConfig,
) -> u32;

pub fn run_substeps(
    state: &mut GlobalState,
    n_substeps: u32,
    config: &Config,
);
\end{verbatim}

\section{Acceptance Tests}\label{sec:sched:tests}

\begin{tabularx}{\textwidth}{l l X}
\toprule
\textbf{ID} & \textbf{Name} & \textbf{Procedure} \\
\midrule
AT-S1 & Disabled passthrough & Set \texttt{enabled: false}; verify
  $n_k = 1$ for all $k$. \\
AT-S2 & Adaptive scaling & Inject high-deformation data; verify
  $n_k > 1$. Inject low-deformation; verify $n_k = n_{\min}$. \\
AT-S3 & Determinism & Run twice with same RD; verify identical $n_k$
  sequences. \\
AT-S4 & Extrinsic invariance & Verify commit index increments exactly
  once per macro-step regardless of $n_k$. \\
AT-S5 & Backward compatibility & Run full S-Basic scenario with
  scheduler disabled; verify identical results to current behavior. \\
\bottomrule
\end{tabularx}


%%% =====================================================================
\part{Engine Extension: Execute Mode}\label{part:execute}

\section{Purpose}\label{sec:exec:purpose}

The execute mode adds a generative operational manifestation to \ISLS{}.
In the existing analytical mode, the engine ingests data and discovers
programs. In execute mode, the engine accepts a program (crystal,
constraint program, or HDAG) as input and executes it through the staged
closure, producing new artifacts.

\section{Formal Model}\label{sec:exec:model}

\begin{definition}[Execute Input]\label{def:execinput}
An execute input is one of:
\begin{enumerate}[nosep,label=(\alph*)]
  \item A \texttt{SemanticCrystal} (the \texttt{constraint\_program}
        field is executed),
  \item A \texttt{ConstraintProgram} (directly),
  \item An HDAG (a \texttt{PersistentGraph} with operator-typed edges).
\end{enumerate}
\end{definition}

\begin{definition}[Execute Semantics]\label{def:execsem}
Given execute input $P$ and initial state $\Sigma_0$ (which may be
empty or pre-populated), the execute mode runs the staged closure:
\[
  \Sigma_K = P(\Sigma_0) \quad \text{under RD}
\]
producing: a final state $\Sigma_K$, zero or more new crystals, a
complete trace, and an execution manifest.
\end{definition}

\begin{definition}[Staged Closure for Execute Mode]\label{def:execstages}
The staged closure in execute mode follows the PCR stages:
\begin{enumerate}[nosep,label=(S\arabic*)]
  \item \textbf{Ingest}: Canonicalize input program $P$.
  \item \textbf{Canon}: Canonicalize initial state $\Sigma_0$.
  \item \textbf{Expand}: Apply macro expansions from the macro registry
        (if $P$ contains macro references).
  \item \textbf{Measure/Embed}: Compute 5D embeddings of the resulting
        graph structure.
  \item \textbf{Solve}: Execute operators in topological order.
  \item \textbf{Gate}: Evaluate all gate obligations (Kairos cascade).
  \item \textbf{Coagula/Select}: If gates pass, crystallize.
  \item \textbf{Emit}: Emit crystals, trace, and manifest.
\end{enumerate}
\end{definition}

\begin{requirement}[Execute-Discover Symmetry]\label{req:symmetry}
The execute mode and the discovery mode \textbf{MUST} share:
\begin{itemize}[nosep]
  \item the same crystal format,
  \item the same evidence chain structure,
  \item the same manifest format,
  \item the same gate cascade,
  \item the same determinism contract.
\end{itemize}
A crystal produced by execute mode \textbf{MUST} be indistinguishable
in format from a crystal produced by discovery mode. The
\texttt{program\_id} field in the manifest distinguishes the provenance.
\end{requirement}

\section{CLI Integration}\label{sec:exec:cli}

\begin{verbatim}
# Execute a constraint program from a crystal file
isls execute --input crystal.json --ticks 50

# Execute an HDAG from a graph file
isls execute --input program.hdag.json --ticks 100 --output results/

# Execute and seal result as a capsule
isls execute --input crystal.json --seal --lock-policy policy.json
\end{verbatim}

\section{Rust API}\label{sec:exec:api}

\begin{verbatim}
// In isls-engine/src/lib.rs (extension)

pub enum ExecuteInput {
    Crystal(SemanticCrystal),
    Program(ConstraintProgram),
    Hdag(PersistentGraph),
}

pub fn execute(
    input: ExecuteInput,
    initial_state: Option<GlobalState>,
    config: &Config,
    registries: &RegistrySet,
    ticks: usize,
) -> Result<(Vec<Option<SemanticCrystal>>, ExecutionManifest)>;
\end{verbatim}

\section{The Discovery-Execute Loop}\label{sec:exec:loop}

The complete PSP substrate enables a self-reinforcing cycle:

\begin{enumerate}[nosep]
  \item \textbf{Discover}: Ingest data, run analytical mode, produce
        crystals encoding discovered programs.
  \item \textbf{Execute}: Take a discovered crystal, run it in execute
        mode on new data, produce new artifacts.
  \item \textbf{Observe}: Ingest the execution artifacts as new
        observations.
  \item \textbf{Re-discover}: Find higher-order patterns in execution
        results.
  \item \textbf{Crystallize}: Produce meta-crystals encoding patterns
        of patterns.
\end{enumerate}

Each cycle produces manifests and evidence chains that are independently
verifiable. The entire history is append-only, content-addressed, and
deterministically replayable.

\begin{theorem}[Loop Determinism]\label{thm:loopdeterminism}
If each iteration of the discovery-execute loop uses a fixed
$\mathrm{RD}$ with fixed registry digests, and all inputs are canonical,
then the entire multi-iteration loop is deterministic: identical initial
observations produce identical meta-crystals at every nesting depth.
\end{theorem}

\begin{proof}
By induction. Each individual run (discover or execute) is deterministic
by Theorem~23.1 of the \ISLS{} v1.0.0 specification. The output of one
run is canonicalized before serving as input to the next. Since
canonicalization is idempotent and all operators are deterministic, the
composition is deterministic at every depth.
\end{proof}


%%% =====================================================================
\part{Dependency Graph and Build Order}\label{part:deps}

\section{Complete Dependency Graph}\label{sec:depgraph}

\begin{verbatim}
isls-types       <-- (all crates)
    ^
isls-registry    <-- (all crates that use operators)
    ^
isls-observe <-- isls-persist <-- isls-extract
                      ^                ^
              isls-consensus <-- isls-carrier
                      ^
              isls-archive <-- isls-morph
                      ^
              isls-manifest  (C13: reads archive + traces)
                      ^
              isls-capsule   (C14: uses manifest + archive + AEAD)
                      ^
              isls-scheduler (C15: wraps engine loop)
                      ^
              isls-engine     (extended: execute mode + scheduler)
                      ^
              isls-harness    (extended: manifest/capsule tests)
                      ^
              isls-cli        (extended: execute + seal/open commands)
\end{verbatim}

\section{Build Order}\label{sec:buildorder}

\begin{enumerate}[nosep]
  \item \texttt{isls-registry} (C12): no new dependencies beyond
        \texttt{isls-types}.
  \item \texttt{isls-manifest} (C13): depends on \texttt{isls-types},
        \texttt{isls-archive}, \texttt{isls-registry}.
  \item \texttt{isls-capsule} (C14): depends on \texttt{isls-types},
        \texttt{isls-manifest}, \texttt{isls-registry}. New external
        dependency: \texttt{aes-gcm} crate for AEAD.
  \item \texttt{isls-scheduler} (C15): depends on \texttt{isls-types},
        \texttt{isls-engine} (or shared types from it).
  \item Engine extensions: modify \texttt{isls-engine} to integrate
        scheduler and execute mode.
  \item CLI extensions: add \texttt{execute}, \texttt{seal},
        \texttt{open} subcommands.
  \item Harness extensions: add tests for C12--C15.
\end{enumerate}

\section{External Dependencies}\label{sec:extdeps}

\begin{tabular}{l l l}
\toprule
\textbf{Crate} & \textbf{Version} & \textbf{Purpose} \\
\midrule
\texttt{aes-gcm} & 0.10 & AEAD encryption for capsules \\
\texttt{hkdf} & 0.12 & Key derivation for capsule session keys \\
\bottomrule
\end{tabular}

All other dependencies are already present in the \ISLS{} workspace.


%%% =====================================================================
\part{Implementation Handoff}\label{part:handoff}

\section{Normative Data Types}\label{sec:handoff:types}

\begin{verbatim}
// New types to add to isls-types/src/lib.rs

pub type Hash256 = [u8; 32];

pub struct SchedulerConfig {
    pub enabled: bool,
    pub n_min: u32,     // >= 1
    pub n_max: u32,     // >= n_min
    pub strategy: String, // "max_pressure" | "weighted" | "fixed"
    pub w_d: f64,       // weight for deformation (if weighted)
    pub w_f: f64,       // weight for friction
    pub w_s: f64,       // weight for shock
}

// Extend RunDescriptor:
pub struct RunDescriptor {
    pub config: Config,
    pub operator_versions: BTreeMap<String, String>,
    pub initial_state_digest: Hash256,
    pub seed: Option<u64>,
    // NEW:
    pub registry_digests: BTreeMap<String, Hash256>,
    pub scheduler: SchedulerConfig,
}
\end{verbatim}

\section{Test Obligations}\label{sec:handoff:tests}

Total new acceptance tests: 21 (AT-R1 through AT-R5, AT-M1 through
AT-M5, AT-C1 through AT-C6, AT-S1 through AT-S5).

All existing 130+ tests \textbf{MUST} continue to pass. The final test
count after extension should be $\ge 151$.

\section{Verification Sequence}\label{sec:handoff:verify}

\begin{verbatim}
# Build
cargo build --workspace --release 2>&1 | grep -c warning
# MUST output: 0

# Test
cargo test --workspace -- --test-threads=1
# MUST output: >= 151 passed, 0 failed

# S-Basic with manifest
rm -rf ~/.isls
cargo run --bin isls --release -- ingest --adapter synthetic --scenario S-Basic
cargo run --bin isls --release -- run --mode shadow --ticks 100
cargo run --bin isls --release -- validate formal
# MUST output: >= 40 crystals, 100% pass rate

# Execute mode
cargo run --bin isls --release -- execute --input ~/.isls/archive/latest.json --ticks 10
# MUST produce crystals and a manifest

# Capsule round-trip
cargo run --bin isls --release -- seal --secret "test" --lock-manifest latest
cargo run --bin isls --release -- open --capsule ~/.isls/capsules/latest.json
# MUST output: "test"
\end{verbatim}

\section{Completion Criteria}\label{sec:handoff:done}

The extension is complete when:
\begin{enumerate}[nosep]
  \item All four crates compile with zero warnings.
  \item All 151+ tests pass.
  \item S-Basic produces crystals with manifests.
  \item Execute mode runs a discovered crystal and produces new crystals.
  \item Seal/open round-trip succeeds.
  \item Scheduler disabled produces identical results to pre-extension
        behavior.
  \item Full HTML report includes manifest and capsule sections.
\end{enumerate}


%%% =====================================================================
\appendix

\section{Invariant Cross-Reference}\label{app:invariants}

\begin{longtable}{l l l}
\toprule
\textbf{New Invariant} & \textbf{Source} & \textbf{ISLS Base Invariant} \\
\midrule
Registry content-addressing & PCR \S5.4 & I2 (Canonical Serialization) \\
Registry drift detection & PCR \S5.4 & I20 (Operator Version Pinning) \\
Manifest content-addressing & PCR \S8.1 & I2 \\
Manifest tamper-evidence & PCR \S8.3 & I3 (Provenance Completeness) \\
Replay pack sufficiency & PCR \S8.2 & I4 (Deterministic Replay) \\
Capsule confidentiality & OLP \S5.2 & new \\
Capsule policy-bound release & OLP \S5.2 & I9 (Threshold-Governed) \\
Capsule tamper-evidence & OLP \S5.2 & I8 (Storage Integrity) \\
Scheduler determinism & NCTE Axiom~3 & I4 \\
Extrinsic discreteness & NCTE Axiom~9 & I10 (Commit Immutability) \\
Execute-discover symmetry & PSP \S I.2 & I2, I4 \\
\bottomrule
\end{longtable}

\section{Symbol Additions}\label{app:symbols}

\begin{tabular}{l l l}
\toprule
\textbf{Symbol} & \textbf{Type} & \textbf{Meaning} \\
\midrule
$\calR$ & Registry & Content-addressed operator/profile/obligation catalog \\
$M$ & Manifest & Execution manifest binding a complete run \\
$T_k$ & TraceEntry & Per-tick trace record \\
$\mathrm{RP}$ & ReplayPack & Self-contained replay bundle \\
$L$ & Lock & $(P_L, \Gamma, \mathrm{RD}_L)$ lock specification \\
$C$ & Capsule & Encrypted evidence-bound artifact \\
$\Pi_C$ & Policy & Capsule release policy \\
$K_M$ & bytes & Master key (external) \\
$K_S$ & bytes & Derived session key \\
$B$ & bytes & Binding string for key derivation \\
$\sigma(\cdot)$ & function & Spiral scheduling function \\
$n_k$ & $\NN_{>0}$ & Sub-step count at macro-step $k$ \\
\bottomrule
\end{tabular}

\vfill
\begin{center}
\rule{0.5\textwidth}{0.4pt}\\[0.5cm]
{\small Generated for \ISLS{} v1.0.0 Extension Architecture.\\
Deterministic, append-only, replay-verified.}
\end{center}

\end{document}
